\section{不動点からの差の比}\label{節：不動点からの差の比}
\subsection{商は長さの比と角度の差}

二点$\alpha,\beta$を基準にして,\,点$z$の位置を
\[
 w=\frac{z-\alpha}{z-\beta}\qquad(z\ne\beta)
\]
で記録してみよう.
絶対値は
\[
 |w|=\frac{|z-\alpha|}{|z-\beta|}
\]
だから,\,二つの基準点からの距離の比を表す.
$z\ne\alpha$なら偏角は,ベクトル$z-\beta$から$z-\alpha$への有向角である.
二つの矢印を同じ始点に平行移動して角度を比較してもよい.
$z=\alpha$では$w=0$となり,偏角は定めない.
\NewFigRatio

この比が一定倍率で変わるなら,\,距離の比と角度の差の変化を同時に追える.
この比を使って,一次分数型の更新を簡単にしよう.

\begin{bookpoint}
複素数の積は\textbf{倍率と回転角}をまとめて記録する.
複素数の商は\textbf{長さの比と角度の差}をまとめて記録する.
\end{bookpoint}

\subsection{二つの不動点を使う}

次の項が
\[
 a_{n+1}=f(a_n),\qquad
 f(z)=\frac{pz+q}{rz+s}
\]
によって決まる場合を考える.
この形の$f$を\textbf{一次分数変換}という.
本節では係数は実数,$r\ne0$,$ps-qr\ne0$とし,\,元の更新で分母が零にならない範囲を扱う.
$r=0$なら$a_{n+1}=Aa_n+B$の形へ戻り,$ps-qr=0$なら定義できる入力で右辺が定数になるため,\,別に調べればよい.

不動点$\alpha$は
\[
 \alpha=\frac{p\alpha+q}{r\alpha+s}
\]
を満たすので,
\[
 r\alpha^2+(s-p)\alpha-q=0
\]
から求まる.
この方程式の根では$r\alpha+s\ne0$である.
もし零なら,\,元の方程式から$p\alpha+q=0$にもなり,$ps-qr\ne0$に反する.

不動点の式を使って引き算すると
\[
 f(z)-\alpha
 =\frac{(p-r\alpha)(z-\alpha)}{rz+s}.
\]
定数項は消えるが,\,まだ分母に$z$が残る.
ここで,\,異なるもう一つの不動点$\beta$があれば
\[
 f(z)-\beta
 =\frac{(p-r\beta)(z-\beta)}{rz+s}.
\]
二式の比を取れば,\,共通の分母が消える.
\[
 \frac{f(z)-\alpha}{f(z)-\beta}
 =\frac{p-r\alpha}{p-r\beta}
  \frac{z-\alpha}{z-\beta}.
\]
つまり
\[
 b_n=\frac{a_n-\alpha}{a_n-\beta},\qquad
 \kappa=\frac{p-r\alpha}{p-r\beta}
\]
と置くと,$b_{n+1}=\kappa b_n$である.
ここで$p-r\alpha,\,p-r\beta$も非零である.
実際,\,不動点の式から
$(p-r\alpha)(r\alpha+s)=ps-qr$となる.

初項が不動点なら定数列として先に扱う.
$a_1\ne\beta$なら,\,上の差の式から$a_n\ne\beta$が続くため,\,比を取れる.
解いた後は
\[
 b_n=b_1\kappa^{n-1},\qquad
 a_n=\frac{\beta b_n-\alpha}{b_n-1}
\]
で元へ戻す.
$b_n=1$になる場合は有限の$a_n$へ戻せず,\,元の更新もその項まで実数列としては定義できない.

\begin{bookpoint}
一つの不動点を引くと定数項が消える.
二つの不動点からの差の比を取ると,\,残った共通の分母も消える.
置き換えは,\,消したい部分に合わせて組み立てられる.
\end{bookpoint}

\noindent\textbf{【例題】}\par\nobreak
$a_1=3$,$a_{n+1}=\frac{2a_n+1}{a_n+2}$を解け.

\noindent\textbf{【方針】}\par\nobreak
不動点1と$-1$からの差の比を取ります.

\noindent\textbf{【解答】}\par\nobreak
\[
 a_n=\frac{2\cdot3^{n-1}+1}{2\cdot3^{n-1}-1}.
\]
\noindent\textbf{【解法】}\par\nobreak
$b_n=\frac{a_n-1}{a_n+1}$とすると
$b_{n+1}=\frac13b_n$,$b_1=\frac12$である.
従って$b_n=\frac1{2\cdot3^{n-1}}$となり,
$a_n=\frac{1+b_n}{1-b_n}$から答えを得る.
分母は正で,\,各項も正なので元の分母$a_n+2$は零にならない.
また$b_n\to0$だから$a_n\to1$となる.
不動点からの距離そのものではなく,\,二つの差の比が縮小していた.

\subsection{不動点が重なる場合}

不動点の方程式が重解$\alpha$を持つと,\,異なる二つの差を作れない.
この場合は
\[
 2r\alpha=p-s,\qquad
 p-r\alpha=r\alpha+s=\ell
\]
となり,$\ell\ne0$である.
$u_n=a_n-\alpha$と置くと
\[
 u_{n+1}=\frac{\ell u_n}{ru_n+\ell}.
\]
初項が不動点の場合を除き,\,逆数を取って
\[
 \frac1{u_{n+1}}=\frac1{u_n}+\frac r\ell
\]
という等差型へ直せる.
二つの固有方向が分かれない場合に足し算が残った\sectionref{節：重解で現れる足し算}と同じように,等比的な変化だけでは捉えられない足し算が残った.

例えば
\[
 a_1=5,\qquad a_{n+1}=\frac{4a_n-9}{a_n-2}
\]
では$\alpha=3$であり,
\[
 \frac1{a_{n+1}-3}=\frac1{a_n-3}+1.
\]
従って$a_n=3+\frac2{2n-1}$となる.
この例では$a_n>3$が保たれ,\,元の分母も非零である.

\subsection{差の比は固有成分の比}

異なる二つの不動点$\alpha,\beta$に対し,\,状態を
\[
 \begin{pmatrix}x_n\\y_n\end{pmatrix}
 =c_n\begin{pmatrix}\alpha\\1\end{pmatrix}
 +d_n\begin{pmatrix}\beta\\1\end{pmatrix}
\]
と分ける. 固有値を$\lambda_\alpha,\lambda_\beta$とすれば
\[
 c_{n+1}=\lambda_\alpha c_n,\qquad
 d_{n+1}=\lambda_\beta d_n.
\]
成分を比較すると
\[
 x_n-\alpha y_n=(\beta-\alpha)d_n,
\]
\[
 x_n-\beta y_n=(\alpha-\beta)c_n.
\]
従って,\,比が定義できる場合には
\[
 \frac{a_n-\alpha}{a_n-\beta}=-\frac{d_n}{c_n}
\]
となり,\,その公比は$\frac{\lambda_\beta}{\lambda_\alpha}$である.
不動点の和$r(\alpha+\beta)=p-s$から
\[
 \lambda_\beta=p-r\alpha,\qquad
 \lambda_\alpha=p-r\beta
\]
なので,\,前の一次分数型の公比と一致する.

本節でも$r\ne0,\,ps-qr\ne0$を仮定する.
このとき$M$は逆行列を持ち,\,固有値も非零である.
不動点を初項とする定数列を除けば,\,必要な成分比を取り続けられる.
元の更新の分母が零にならない条件は別に確かめる.

実係数で不動点が共役な非実数なら,\,固有値も共役である.
両者は絶対値が等しいため,\,\textbf{比を取ると共通の拡大率が消えて回転だけが残る}.
\sectionref{節：隣接3項間漸化式と回転}では回転に拡大・縮小が伴った. 次節で扱う一次分数型の円周上の座標では,成分比を取ることで共通の倍率が消える.

不動点が重なる場合は,\,行列も二つの固有方向に分けられない.
$r\ne0$だから$M$は単位行列の定数倍ではなく,\,固有方向は一本になる.
その寄与の蓄積を比で読み直すと,\,逆数を取った等差型が現れる.
隣接3項間漸化式の重解と,一次分数型の重なった不動点は,同じ行列の性質から生じている.
